\documentclass[aspectratio=169,11pt]{beamer}

% ── Session 1: "The Wrench in the Works" ──────────────────────────────────
% Format-locking deck for the Great Wall pilot. Lively on purpose: humor,
% memes, big type. Memes are loaded via \IfFileExists guards so a missing
% image never breaks the build — drop real PNGs into ../assets/ (see
% ../assets/README.md) and they appear automatically.
% Shared brand theme + \gwtitle/\punch/\memeslot live in gwtheme.tex.

\usepackage[utf8]{inputenc}
\usepackage[T1]{fontenc}
\usepackage{lmodern}
\usepackage{../style/greatwall}   % shared brand palette + boxes + marks
\usepackage{booktabs}
\usepackage{array}                % \arraybackslash for p-column alignment
\usepackage{colortbl}             % \rowcolor in tables

\graphicspath{{./}{../assets/}}
\def\gwsessionno{1}
% gwtheme.tex — shared Beamer theme for the Great Wall pilot-course decks.
% \input this from each slides/sessionN.tex *after* loading greatwall.sty and
% *after* defining \gwsessionno (the session number, e.g. \def\gwsessionno{2}).
% Provides the brand look plus three helpers: \gwtitle, \punch, \memeslot.
%
% Depends on: beamer, greatwall.sty (palette + boxes + marks), tikz.

\providecommand{\gwsessionno}{?}

% ── Minimal custom theme (no external font deps; texlive-full only) ────────
\usefonttheme{professionalfonts}
\setbeamertemplate{navigation symbols}{}
\setbeamercolor{normal text}{fg=darktext, bg=white}
\setbeamercolor{structure}{fg=accent}
\setbeamercolor{frametitle}{fg=white, bg=accent}
\setbeamercolor{title}{fg=white}
\setbeamercolor{block title}{fg=white, bg=accentdark}
\setbeamercolor{block body}{fg=darktext, bg=accentlight}
\setbeamercolor{itemize item}{fg=accent}
\setbeamercolor{itemize subitem}{fg=btccol}
\setbeamerfont{frametitle}{series=\bfseries}
\setbeamerfont{title}{series=\bfseries, size=\Huge}
\setbeamertemplate{itemize item}{\textbf{--}}

% Frametitle as a solid green bar
\setbeamertemplate{frametitle}{%
  \nointerlineskip
  \begin{beamercolorbox}[wd=\paperwidth,ht=2.6ex,dp=1.2ex,leftskip=0.4cm,rightskip=0.4cm]{frametitle}%
    \usebeamerfont{frametitle}\insertframetitle%
  \end{beamercolorbox}%
}

% Footline: course | session | frame number
\setbeamercolor{gwfootline}{fg=graytext, bg=white}
\setbeamertemplate{footline}{%
  \leavevmode%
  \hbox{%
  \begin{beamercolorbox}[wd=.5\paperwidth,ht=2.4ex,dp=1ex,leftskip=0.4cm]{gwfootline}%
    \tiny Great Wall --- Pilot Course%
  \end{beamercolorbox}%
  \begin{beamercolorbox}[wd=.5\paperwidth,ht=2.4ex,dp=1ex,rightskip=0.4cm]{gwfootline}%
    \hfill\tiny Session \gwsessionno{} \;\textbullet\; \insertframenumber%
  \end{beamercolorbox}}%
}

% ── Helpers ───────────────────────────────────────────────────────────────
% Big punchy centered statement frame (full-bleed green).
\newcommand{\punch}[1]{%
  \begin{frame}[plain]
  \begin{tikzpicture}[remember picture,overlay]
    \fill[accent] (current page.south west) rectangle (current page.north east);
  \end{tikzpicture}
  \begin{center}\vfill
  {\color{white}\Huge\bfseries #1}\par
  \vfill\end{center}
  \end{frame}}

% Meme slot — shows the image if present, else a labelled placeholder box.
% #1 = filename in ../assets/, #2 = caption/placeholder text.
\newcommand{\memeslot}[2]{%
  \IfFileExists{../assets/#1}{%
    \includegraphics[width=\linewidth,height=0.62\textheight,keepaspectratio]{#1}%
  }{%
    \begin{tcolorbox}[colback=accentlight,colframe=graytext,sharp corners,
      halign=center, valign=center, height=0.55\textheight]
      \color{graytext}\itshape [meme slot: \texttt{#1}]\\[4pt]\small #2
    \end{tcolorbox}%
  }%
}

% Branded title frame. #1 = title, #2 = subtitle.
\newcommand{\gwtitle}[2]{%
  \begin{frame}[plain]
    \begin{tikzpicture}[remember picture,overlay]
      \fill[accent] (current page.south west) rectangle (current page.north east);
      \foreach \i in {1,...,8}{
        \fill[white,opacity=0.08] ([xshift=-\i cm,yshift=\i cm]current page.south east)
          rectangle ++(0.7cm,0.7cm);
      }
    \end{tikzpicture}
    \vfill
    \begin{center}
    {\color{white}\Huge\bfseries #1\par}
    \vspace{6pt}
    {\color{accentlight}\large #2\par}
    \vspace{18pt}
    {\color{accentlight}\small Great Wall \;\textbullet\; Session \gwsessionno{} of 5 \;\textbullet\; hands-on, on Discord\par}
    \end{center}
    \vfill
  \end{frame}}


\begin{document}

\gwtitle{The Wrench in the Works}{Why self-custody isn't enough --- and what we'll do about it}

% ── Check-in ──────────────────────────────────────────────────────────────
\begin{frame}{Before we start: a one-word check-in}
  \begin{columns}[T]
    \begin{column}{0.58\textwidth}
      Drop \textbf{one word} in \texttt{\#feedback} for how you arrived today.
      \vspace{8pt}

      And a quick poll --- react in \texttt{\#live-chat}:
      \begin{itemize}
        \item[\textbf{\color{accent}(A)}] ``I could lose my crypto to a hacker.''
        \item[\textbf{\color{btccol}(B)}] ``I could lose my crypto to a \emph{person} standing next to me.''
      \end{itemize}
      \vspace{6pt}
      Hold that thought. By the end of today, \textbf{\color{btccol}(B)} is the one we care about.
    \end{column}
    \begin{column}{0.40\textwidth}
      \begin{block}{House rules}
        \footnotesize
        Throwaway setups + test entropy \textbf{only}. Never a real seed, never
        real funds --- not even to show off.
      \end{block}
    \end{column}
  \end{columns}
  \note{Keep this to 3 minutes. The poll seeds the whole session's tension.}
\end{frame}

% ── The hook: the headline ────────────────────────────────────────────────
\begin{frame}{The headline that started this}
  \begin{columns}[T]
    \begin{column}{0.52\textwidth}
      \begin{block}{X, 12 Feb 2026}
        \footnotesize\itshape
        ``ATTACKS ON BITCOINERS IN FRANCE ACCELERATING --- Out of 47 documented
        attacks since 2017, 38 happened in the last 13 months.''
      \end{block}
      \vspace{4pt}
      The top reply, endorsed as \emph{``golden advice''}:
      \begin{block}{}
        \footnotesize\itshape ``No one needs to know you own Bitcoin.''
      \end{block}
    \end{column}
    \begin{column}{0.46\textwidth}
      \memeslot{wrench-attack-meme.png}{the classic \$5 wrench xkcd vibe ---
      ``\$5 wrench'' \(>\) any amount of crypto}
    \end{column}
  \end{columns}
\end{frame}

% ── Define wrench attack ──────────────────────────────────────────────────
\begin{frame}{The \texorpdfstring{\wrench}{wrench attack}}
  \begin{keyidea}[The threat we actually can't ignore]
    A \textbf{wrench attack} (a.k.a.\ ``\$5 wrench attack''): forget hacking the
    cryptography --- just physically coerce the human who holds the keys.
    \emph{``Transfer it, or else.''}
  \end{keyidea}
  \vspace{4pt}
  Your hardware wallet's secure element, your 24 words on titanium, your fancy
  passphrase --- none of it helps when the attacker is in the room and patient.
  \vspace{6pt}
  \begin{center}\small\color{graytext}
  Tracked in the real world: \texttt{jlopp/physical-bitcoin-attacks}.
  \end{center}
\end{frame}

% ── The contradiction (punch) ─────────────────────────────────────────────
\punch{If everyone is a needle,\\[6pt] there is no hay.}

\begin{frame}{Obscurity vs.\ adoption: pick one}
  The consensus defense is \textbf{obscurity}: ``tell no one.'' But:
  \begin{itemize}
    \item \textbf{It fights adoption.} Safety-by-blending-in cannot scale to
      everyone --- if all Bitcoiners are needles, there's no haystack left.
    \item \textbf{It violates Kerckhoffs (1883).} Security must not rest on the
      \emph{secrecy of the method}. Bitcoin is open; an educated attacker knows
      what you know.
    \item \textbf{Decoy tricks backfire.} The attacker can't tell a decoy wallet
      from the real one --- so they just\ldots keep going.
  \end{itemize}
  \vspace{4pt}
  \begin{block}{The community endorses obscurity \emph{in the same threads} where it documents obscurity failing.}
  \end{block}
\end{frame}

% ── Skeptic seat intro ────────────────────────────────────────────────────
\begin{frame}{Interaction: the skeptic's seat}
  \begin{columns}[T]
    \begin{column}{0.6\textwidth}
      One volunteer per session has a job: \textbf{attack the claim.}
      \vspace{6pt}

      Today's claim to attack:
      \begin{block}{}
        ``Obscurity is a \emph{cost}, not a defense.''
      \end{block}
      \vspace{6pt}
      Everyone else: back them up or knock them down in \texttt{\#live-chat}.
    \end{column}
    \begin{column}{0.38\textwidth}
      \memeslot{skeptic-seat.png}{``but couldn't the attacker just\ldots ?'' ---
      yes. that's the point of this seat.}
    \end{column}
  \end{columns}
  \note{Assign the skeptic now; they hold the role all session.}
\end{frame}

% ── The data: security tax ────────────────────────────────────────────────
\begin{frame}{Let's put a price on it}
  Insurance markets price ``portable-valuable-gets-stolen'' every day (jewelry,
  watches). Borrow their math for Bitcoin whales:
  \begin{keyidea}[The security tax]
    At the 2026 pace, the insurance-equivalent annual security cost for individual
    BTC whales has entered the \textbf{gold-storage range} ($\sim$0.5--1.5\%/yr) ---
    a cost the market currently prices at \textbf{zero}.
  \end{keyidea}
  And because attacks grow faster than the whale population, the
  \emph{forward-looking} discount is several times the current annual number.
  \begin{center}\footnotesize\color{graytext}
  Full derivation: \texttt{justification-and-economics/JUSTIFICATION}.
  \end{center}
\end{frame}

\begin{frame}{Attacks are outgrowing the targets}
  \begin{columns}[T]
    \begin{column}{0.56\textwidth}
      \begin{tikzpicture}[scale=1]
        \draw[->,thick,graytext] (0,0) -- (5.4,0) node[right]{\scriptsize year};
        \draw[->,thick,graytext] (0,0) -- (0,3.4) node[above]{\scriptsize rate};
        % attack rate rising
        \draw[attackcol,very thick]
          (0.2,0.4) -- (1.2,0.7) -- (2.2,1.4) -- (3.2,1.6) -- (4.2,2.6) -- (5.0,3.1);
        \node[attackcol,right] at (3.0,2.9) {\scriptsize attacks/100k};
        % whale pop rising slower
        \draw[accent,very thick,dashed]
          (0.2,0.6) -- (1.2,0.9) -- (2.2,1.1) -- (3.2,1.3) -- (4.2,1.5) -- (5.0,1.7);
        \node[accent,right] at (3.4,1.2) {\scriptsize whales};
      \end{tikzpicture}
    \end{column}
    \begin{column}{0.42\textwidth}
      The red line (attacks per capita) is pulling away from the green
      (population). The \emph{per-person} risk is rising.
      \vspace{6pt}
      \begin{block}{We are early in the trend}
        \footnotesize Targeting gets easier and criminal skill diffuses as
        adoption grows. This converges toward ordinary-robbery economics.
      \end{block}
    \end{column}
  \end{columns}
  \gwfeedback{Does this read as evidence or as a pitch? Genuinely asking ---
  it's the hook the course hangs on.}
\end{frame}

% ── The four-property test ────────────────────────────────────────────────
\begin{frame}{The four-property test}
  Four properties an adequate method needs \emph{at once}. Every method on the
  market drops one --- and the dropped one is the one that gets people hurt.
  \vspace{4pt}
  \begin{center}\footnotesize
  \renewcommand{\arraystretch}{1.35}
  \begin{tabular}{@{}l c c c c >{\raggedright\arraybackslash}p{0.30\textwidth}@{}}
    \toprule
    \textbf{Method} & \textbf{Devi.} & \textbf{Indiv.} & \textbf{Anti-obsc.} & \textbf{Coerc.} & \textbf{Fails on} \\
    \midrule
    \rowcolor{accentlight}\textbf{TKBA} (Great Wall) & \gwcheck & \gwcheck & \gwcheck & \gwcheck & only: \emph{forgetting} (we solve it) \\
    Physical (vault, Shamir) & \gwcross & \gwcheck & \gwcheck & \gwcheck & cost + ``glorified gold'' \\
    Shared (Coinbase, Casa) & \gwcheck & \gwcross & \gwcheck & \gwcheck & not your keys\ldots \\
    Obscurity (decoy, denial) & \gwcheck & \gwcheck & \gwcross & \gwskull & \emph{backfires} $\to$ more violence \\
    Vanilla (Ledger, Electrum) & \gwcheck & \gwcheck & \gwcheck & \gwcross & wide open to the wrench \\
    \bottomrule
  \end{tabular}
  \end{center}
  \begin{center}\scriptsize\color{graytext}
  \gwcheck~has it \quad \gwcross~lacks it \quad \gwskull~\emph{negative} (worse than lacking) \quad\textbf{\color{accentdark}Only the top row holds all four.}
  \end{center}
  \gwfeedback{The obscurity row's \gwskull{} (not a \gwcross) is the subtle bit ---
  does ``negative resistance'' land? We unpack it next.}
\end{frame}

% ── Two cases: Balland vs RATIRL ──────────────────────────────────────────
\begin{frame}{Two cases that bracket the problem}
  \begin{columns}[T]
    \begin{column}{0.48\textwidth}
      \begin{block}{Force beats a fortress --- David Balland}
        \footnotesize
        Co-founder of \textbf{Ledger}. Wealthy, expert, security-aware --- about as
        well-defended as an individual gets.
        \vspace{3pt}
        Jan 2025: kidnapped $\sim$2 days, \textbf{mutilated}, ransom extorted;
        freed only by police.
      \end{block}
      \scriptsize\color{graytext} Abundant physical defense didn't prevent the
      attempt \emph{or} the harm.
    \end{column}
    \begin{column}{0.48\textwidth}
      \begin{block}{Denial gets you tortured --- RATIRL}
        \footnotesize
        Swedish \emph{LoL} streamer. Armed intruders, \textbf{tortured for hours}
        over crypto he says he didn't have.
        \vspace{3pt}
        They didn't believe him; left with his \textbf{Pok\'emon cards}.
      \end{block}
      \scriptsize\color{graytext} By his own account the ordeal was
      \emph{prolonged} --- the tell that the denial was \emph{pre-disbelieved}.
    \end{column}
  \end{columns}
  \vspace{6pt}
  \begin{center}\small
  \textbf{\color{accentdark}Bodyguards can't fix it (Balland). Talking can't fix
  it (RATIRL).}
  \end{center}
\end{frame}

% ── The rhetoric trilemma ─────────────────────────────────────────────────
\begin{frame}{The rhetoric trilemma}
  Under coercion, everything you can \emph{do} is a \emph{speech act} the
  \textbf{attacker} judges. Only three moves exist --- all lose:
  \vspace{2pt}
  \begin{center}
  \begin{tikzpicture}[font=\footnotesize, node distance=6mm,
    m/.style={draw=accentdark, rounded corners=2pt, fill=accentlight,
              text width=2.7cm, align=center, minimum height=1.05cm, inner sep=3pt},
    f/.style={draw=attackcol, rounded corners=2pt, fill=white, text=attackcol,
              text width=2.7cm, align=center, minimum height=1.05cm, inner sep=3pt},
    ar/.style={-latex, thick, attackcol}]
    \node[m] (c) {\textbf{Comply}\\\scriptsize hand over the key};
    \node[m, right=of c] (d) {\textbf{Deny}\\\scriptsize ``I have nothing''};
    \node[m, right=of d] (k) {\textbf{Decoy}\\\scriptsize a small / fake wallet};
    \node[f, below=of c] (cf) {robbed\\\scriptsize (and maybe not done)};
    \node[f, below=of d] (df) {presumed \emph{denial-script} $\to$ escalate};
    \node[f, below=of k] (kf) {presumed \emph{the decoy} $\to$ escalate};
    \draw[ar] (c)--(cf); \draw[ar] (d)--(df); \draw[ar] (k)--(kf);
  \end{tikzpicture}
  \end{center}
  \vspace{-2pt}
  \begin{center}\footnotesize
  All three are \emph{unverifiable claims}; obscurity culture pre-discredits the
  last two. Honest ``have-nothing'' and stonewaller look \textbf{identical} ---
  so the rational next step is \emph{more} violence.
  \end{center}
\end{frame}

% ── How Great Wall escapes it ─────────────────────────────────────────────
\begin{frame}{How Great Wall breaks the trilemma}
  \begin{keyidea}[Swap rhetoric for a verifiable constraint]
    ``I \emph{cannot} produce it within your window'' stops being a \emph{claim to
    be believed} and becomes a \textbf{property of a public protocol the attacker
    can check for themselves}.
  \end{keyidea}
  \vspace{2pt}
  \begin{itemize}\small
    \item \textbf{Balland's lesson:} put the barrier at the \emph{protocol} layer
      (nothing to hand over), not the bodyguard layer that force overruns.
    \item \textbf{RATIRL's lesson:} the victim's signal must be \emph{verifiable},
      not rhetorical.
    \item \textbf{Anti-Obscurity is the trick:} the one honest signal is credible
      \emph{because} the attacker knows the mechanism --- exactly the property
      obscurity throws away.
  \end{itemize}
\end{frame}

% ── The reveal ────────────────────────────────────────────────────────────
\punch{What if the secret were\\[6pt] something you literally\\[6pt] \emph{cannot} hand over?}

\begin{frame}{The Four Properties --- all at once}
  Great Wall claims all four \emph{simultaneously}:
  \begin{enumerate}
    \item \textbf{Knowledge-based} --- it's entirely in your memory.
    \item \textbf{Individual} --- you depend on no one else.
    \item \textbf{Non-obscure} --- the method is public; hiding nothing.
    \item \textbf{Coercion-resistant} --- threats don't get the secret.
  \end{enumerate}
  \vspace{6pt}
  \begin{keyidea}[The whole thing in one sentence]
    It's all in your head (1), in nobody else's (2), the attacker \emph{knows}
    that (3), and is nevertheless unable to rob it (4).
  \end{keyidea}
\end{frame}

\begin{frame}{``That sounds like having your cake and eating it.''}
  \begin{columns}[T]
    \begin{column}{0.58\textwidth}
      Correct. It's possible only because the secret is \textbf{tacit}
      (you can recognize it but not dictate it) and deploying it is gated by an
      \textbf{unavoidably slow computation}.
      \vspace{8pt}

      We don't ask you to believe that today. We'll \emph{build} it:
      \begin{itemize}\small
        \item S2 --- secrets you can't tell (tacit knowledge)
        \item S3 --- drawing inside the labyrinth (the fractal)
        \item S4 --- building the wall (Argon2)
        \item S5 --- living with it (train / accelerate / inherit)
      \end{itemize}
    \end{column}
    \begin{column}{0.40\textwidth}
      \memeslot{cake-eat-it.png}{``have your cake and eat it'' --- except this
      time the math actually allows it}
    \end{column}
  \end{columns}
\end{frame}

% ── Roadmap ───────────────────────────────────────────────────────────────
\begin{frame}{Where we're going (5 sessions)}
  \begin{center}
  \begin{tikzpicture}[node distance=4mm and 7mm, font=\footnotesize,
    s/.style={draw=accentdark, fill=accentlight, rounded corners=2pt,
              text width=2.0cm, align=center, minimum height=1.5cm},
    a/.style={-latex,thick,accent}]
    \node[s,fill=accent,text=white] (s1) {\textbf{1}\\The Problem\\\textit{(today)}};
    \node[s,right=of s1] (s2) {\textbf{2}\\Tacit\\knowledge};
    \node[s,right=of s2] (s3) {\textbf{3}\\Fractal\\encoding};
    \node[s,right=of s3] (s4) {\textbf{4}\\The Argon2\\wall};
    \node[s,right=of s4] (s5) {\textbf{5}\\Live setup\\+ ecosystem};
    \draw[a] (s1)--(s2); \draw[a] (s2)--(s3); \draw[a] (s3)--(s4); \draw[a] (s4)--(s5);
  \end{tikzpicture}
  \end{center}
  \vspace{4pt}
  \begin{center}\itshape\color{graytext}
  Every session ends at a keyboard. Today's keyboard moment is your homework.
  \end{center}
\end{frame}

% ── Hands-on homework ─────────────────────────────────────────────────────
\begin{frame}{Hands-on: get lost in a fractal}
  \begin{tryit}[Homework into Session 2]
    \begin{enumerate}\small
      \item Clone \texttt{great-wall-core}; build the Rust engine
        (\texttt{cargo build --release}).
      \item Launch the viewer: \texttt{cd burning\_ship \&\& python3 viewer.py}
      \item \textbf{Just play.} Pan, zoom, cycle color schemes (\texttt{1}--\texttt{5}).
        No goal. Get a feel for how strange and infinitely-detailed the Burning
        Ship is.
      \item Post your favorite view in \texttt{\#show-your-fractal}.
    \end{enumerate}
  \end{tryit}
  Stuck on the build? That's what \texttt{\#hands-on-help} (and the optional
  ``Session~0'' clinic) is for. \emph{Don't} burn your own time silently ---
  install pain is data we want.
\end{frame}

% ── Reflect ───────────────────────────────────────────────────────────────
\begin{frame}{Reflect (one paragraph)}
  Write your own \textbf{threat profile}:
  \begin{itemize}
    \item Who could plausibly coerce you?
    \item For how long could they realistically hold you?
    \item Which conventional defense are you relying on right now --- and which
      hole does it have?
  \end{itemize}
  \vspace{6pt}
  We build directly on this next week. There are no wrong answers --- only
  un-examined ones.
\end{frame}

% ── Check-out ─────────────────────────────────────────────────────────────
\begin{frame}{Check-out --- 60 seconds}
  \begin{columns}[T]
    \begin{column}{0.6\textwidth}
      \begin{block}{Two quick things in \texttt{\#feedback}}
        \begin{enumerate}\small
          \item Your \textbf{muddiest point} today, in one line.
          \item The per-session form (link pinned) --- 3 minutes, please do it
            \emph{now} while it's fresh.
        \end{enumerate}
      \end{block}
      \vspace{4pt}
      \footnotesize Especially: did the security-tax framing convince you, or
      feel like marketing? Be brutal. That's the gift.
    \end{column}
    \begin{column}{0.38\textwidth}
      \begin{center}
      \gwflagicon\\[4pt]
      {\footnotesize\color{attackcol}\itshape Your confusion is the curriculum.}
      \end{center}
    \end{column}
  \end{columns}
\end{frame}

% ── Closing ───────────────────────────────────────────────────────────────
\punch{See you in the labyrinth.}

\end{document}
